% -----------------------------------------------------------
\section{NT Extended Scripture Studies}
% -----------------------------------------------------------
	\label{EXSS-NT}
	
	% ----------------------
	\subsection{The Synoptic Problem}
	% ----------------------
		\label{EXSS-NT-SYNOPTIC}
		
		There are at least two aspects to the synoptic problem:
		
			\begin{compactitem}
				\item The proposed dependence of Matthew and Luke on Mark for the material that all three have in common;
				\item The proposed common-source ``Q'' as the source for the material that Matthew and Luke have in common but that does not appear in Mark.
			\end{compactitem}
			
		As ABD 8529 puts it, ``The three Synoptic Gospels are thus in some relationship with each other, and the problem of determining the nature of that relationship is known as the synoptic problem.''
			
		See the following reference sources for more information:
		
			\begin{compactdesc}
				\item[ABD] Q (Gospel Source); Source Criticism (NT); Synoptic Problem; Two-Gospel Hypothesis; Two-Source Hypothesis
			\end{compactdesc}
	
	% ----------------------
	\subsection{Matthew 5:3--12: On the Beatitudes in the Sermon on the Mount}
	% ----------------------
		\label{EXSS-NT-BEATITUDE}
	
		% -----
		\subsubsection{NIVAC NT 01 PDF 315--319}
		% -----
		
			Dallas Willard's influential book \textit{The Divine Conspiracy} is a forceful development of Christian discipleship, concentrating on the Gospel of Matthew and the SM in particular as his biblical foundation. In his discussion of the Beatitudes, he recounts how a woman told him that her son had dropped his Christian identification and left the church because of the Beatitudes. This son, a strong, intelligent, military person, had had an unhappy experience: 
		
				\begin{quote}
					As often happens, he had been told that the Beatitudes---with its list of the poor and the sad, the weak and the mild---were a picture of the ideal Christian. He explained to his mother very simply: ``That's not me. I can never be like that.'' 
				\end{quote}

			I understand this young man's experience. In my teens I also renounced my Christian identity and church for similar reasons. I clearly remember sitting in youth group meetings where the characteristics of the Beatitudes were held up as ideals for us to emulate. I remember snickering in the back row with my buddies as the youth leaders cajoled us to cease being cocky and macho and become meek and mild. The four of us were three-sport athletes in high school, and the picture of the Christian life that was held up for us from the Beatitudes seemed lamely pathetic. As I think back, our youthful cockiness and machismo were probably just as pathetic, but the Christian life painted by that church had nothing to offer us as a viable, robust alternative. 

			Not too many years after ruling out the Beatitudes for real life, I sat under the brilliant stars in a jungle in Vietnam and their significance overwhelmed me. I was a member of a cocky airborne infantry combat battalion. We were a well-trained, exceedingly efficient war machine. One night as I sat on guard duty after one especially ravaging battle, I experienced the reality of what Jesus addressed in the Beatitudes. I had killed gleefully that day. I had ripped the life from other young men without a twinge of conscience. I saw the bodies of my nineteen-and twentyyear-old squad members ravaged by other young men who were our hated enemies, yet probably none of us on either side could really offer any adequate explanation for our animosity. 

			That night I experienced brokenness. I became poor in spirit as I recognized the depth of my depravity and shuddered as I considered the possibility of my fate before God, if he existed. I mourned at the evil in me and at the evil that I saw emerge so quickly in all of us. For the first time in my young life, I understood that I was not the invincible captain of my ship. I could be killed at any moment. So from that very night I began to realize that there was indeed a very different way to live. I did not articulate it that night in these words, but meekness, righteousness, mercy, purity, and peacemaking all became so much more clearly preferable than the way that I had been pursuing significance and success. 

			I now realize I was experiencing the beginnings of the \textit{pronouncement} aspects of the Beatitudes. I saw for the first time the horror of my life as a human apart from God. I desperately needed something, but what it was, I had no clue. I experienced the condemnation of my old cockiness and self-sufficiency, and above all, the condemnation of my arrogant abuse of people in my quest to satisfy my own lusts. This transition in my life readied me and enabled me to accept Jesus' invitation to the life of the kingdom of heaven two years later... 

			In the same way that the Beatitudes express the blessedness that comes to the crowds from the convicting work of the Holy Spirit, they also express the blessedness that comes from the renewing work of the Holy Spirit in the life of a disciple... 

			The individual characteristics of the Beatitudes are not self-produced, nor can we simply learn or emulate them in an attempt to bring them about in our lives. They are products of a life energized by the Spirit of God. They are, like the listing Paul gives in \hyperref[GALATIANS-5-22]{Galatians 5:22--23}, the fruit of the Spirit. They are a wholistic view of what the Spirit will produce in the life of a disciple of Jesus who is walking in his ways and is being transformed into his image. 

			So it does help tremendously to study the Beatitudes, because they reveal the values of the kingdom of heaven. As in any study of Scripture, they show us God's ways in distinction from the world's ways and help us to know the right path. But the wonderful truth behind the study of the Beatitudes and our obedience to their truth as the Word of God is that the characteristics of the Beatitudes are ultimately produced by the Spirit of God. 

			The young man mentioned earlier had turned away from Christianity because of the perception he had of the Christian life displayed in the Beatitudes. Two things are probably at play here. (1) Like me, he is a product of a world that glories in self, in personal and institutional strength and bravado. The Beatitudes turn those values on end to show what humanity rightly related to God and to each other will be through the arrival of the kingdom of God. 

			(2) He quite likely has seen faulty applications of the Beatitudes in his church experiences. Throughout church history the Beatitudes have been subjected to faulty interpretation that has led to extremes. Many years later I came to understand that the church I left in my teens was theologically liberal. They really did believe that the Beatitudes were an expression of the ideal life that humans need to pursue to find God. They rightly recognized that the Beatitudes are an ideal statement of the Christian life, but they mistakenly thought that they could do it on their own. 

			Some groups have contended that these characteristics are not for today's rough-and-tumble, sinwracked world, but await some far-off future kingdom. Still others have based entire personal and ecclesiastical practices of pacifism and nonviolence on the Beatitudes' overriding centrality for a present-day theological system. No wonder the young man did not see what Jesus meant to communicate. The Beatitudes are neither a means of entering nor of advancing in the kingdom. They are expressions of Spirit-produced kingdom life, revealing to the entire world that a transformation of creation is beginning in Jesus' disciples. That is why we are blessed. 
			
	% ----------------------
	\subsection{Matthew 10:39: On Losing One's Self}
	% ----------------------
		\label{EXSS-NT-MATTHEW-10-39-SELF}
		
		% -----
		\subsubsection{WBC 33a PDF 370}
		% -----
		
			Taking up one's cross refers not to the personal problems or difficulties of life that one must bear, as it is sometimes used in common parlance, but to a radical obedience that entails self-denial and, indeed, a dying to self. To take up one's cross is to follow in the footsteps of Jesus, who is the model of such radical obedience and self-denial (cf. 4:1-11). Thus in a real sense v 39 is a kind of exegesis of v 38. ``Finding one's life'' (\gr{εὑρὼν τὴν ψυχὴν}) in v 39a refers obviously to the affirmation of life on one's own terms within one's self-centered framework apart from allegiance and discipleship to Jesus. ``Find'' here means ``to obtain'' or ``find for oneself'' (BAGD, 325 [BDAG 363b meaning \#3]). This person will in the end lose his or her life, i.e., will not inherit the life of the future kingdom. On the contrary, the one who loses his or her life \gr{ἕνεκεν ἐμοῦ}, ``for my sake,'' i.e., in the self-denial entailed in the taking up of one's cross, which might mean even martyrdom (vv 21, 28; cf. Rev. 12:11), is ultimately the one who will ``find'' it, i.e., in a degree of fulfillment in the present age but preeminently in the age to come (cf. John 12:25; for a Jewish parallel, see \textit{Abot R Nat.} II[2]. 36a). For the sake of Jesus, and thus the gospel, the disciple is called to follow after Jesus in unqualified obedience to the will of God, even to the point of death itself, which becomes for the disciple the entry into life. \gr{ψυχή} can mean either ``soul'' (cf. 10:28) or ``life'' in its more inclusive sense (e.g., 6:25). Since literal martyrdom is not required to ``find life'' (though it may be the lot of some), losing life and finding life are here taken in the extended sense of meaningful existence, fulfillment, purpose, or identity. There is no real life in this sense apart from relationship to Jesus.
			
	% ----------------------
	\subsection{Luke 8:5--13: On Prayer}
	% ----------------------
		\label{EXSS-NT-LUKE-8-PRAYER}
		
		% -----
		\subsubsection{NIVAC NT 01 PDF 3153--3154, 3158--3160}
		% -----
		
			To understand this parable, we must appreciate certain aspects of first-century cultural expectation. (1) Food was not as readily available as it is today; the era of twenty-four hour convenience stores is new. Bread was baked each day to meet daily needs. (2) The culture held hospitality in high regard, almost as a duty. A visitor was to be welcomed and cared for, regardless of the hour of his arrival. Here, then, we have a dilemma: a late evening guest, but no food. What would the host do? Also behind the story stands the reality that most ancient homes had only one room. To approach a neighbor meant risking waking the family. How bold would the host be? 

			Jesus poses the dilemma and then asks, ``Which of you would have the nerve to go wake up his friend (and possibly his family as well!) in the middle of the night to ask for bread?'' The request for three loaves would meet the need, provided his neighbor still had some leftovers from the day. The answer shows the tension. He replies that he cannot be bothered, since the children are in bed and he is unwilling to risk rousing them! Anyone who has put children to bed understands how the neighbor feels. In fact, the natural quality of the reply is part of the charm of this parable. Jesus then comes to his main point. The neighbor does respond, not because he is a friend but because the petitioner has ``boldness'' (\textit{anaideian}). This Greek term is difficult to render in English. It combines two qualities: boldness and shamelessness. This man has nerve to make such a request. In recognition of his boldness, the neighbor honors the request (cf. \hyperref[HEBREWS-10-19]{Heb. 10:19--22}). 

			The qualification of these requests as applying to basic spiritual needs comes from the concluding verses of the unit (vv. 9--10), where Jesus issues a call to ``ask,'' ``seek,'' and ``knock.'' The promise of an open door lies before the disciple, who need not be too shy to ask. Then Jesus gives another illustration. If a child asks for a basic meal, fish, or eggs, will the father give him something dangerous instead, like a snake or scorpion? What father would do that? None. So Jesus concludes, ``If you then, though you are evil, know how to give good gifts to your children, how much more will your Father in heaven give the Holy Spirit to those who ask him!''... 

			God is not a cosmic grouch. He loves and cares for us, desiring what is best for us. He shows the extent of his commitment to us in the sacrifice of Jesus (\hyperref[ROMANS-8-30]{Rom. 8:30--39}). Nothing can separate us from a close connection with him but ourselves and our failure to accept the care, provision, and access that comes from his hands. A relationship takes more than words claiming his presence; it must be nurtured through time and effort. That is why prayer is so crucial to the Christian walk. If Jesus took the time to teach us how to pray and urged us to do so with boldness, then our calling is to make time for it. 

			The pursuit of spiritual discipline is receiving renewed attention today. This is only right. Our world is so hectic that often all we think we have time for is a quick ``check in'' with the Father above. In fact, however, the more hectic the pace, the more we need him. In Jesus' life, Luke notes him praying in the moments when things were getting the most hectic. He prayed when he was baptized, before he chose the Twelve, after Peter's confession at the Transfiguration as he prepared to face Jerusalem, and when he was in the Garden, even as the guards approached to take him prisoner---not to mention on the cross itself. When events surrounded him, Jesus circled the wagons around the courts of God. The disciples must have sensed this, because this is one of the few places where they make a specific request to be taught something. The application is simple: Make time to see his face. 

			Another point emerges in the juxtaposition of intimacy with respect for God's uniqueness. God is not a ``buddy,'' he is both Father and Lord. The petitioner who recognizes his uniqueness appreciates the privilege of entering his presence. Kneeling before the God of the universe requires humility. In our culture, where respect for elders often has reversed itself into a worshiping of youth, it is important to appreciate that a relationship with God is not a peer relationship. He sits enthroned in the heavens above. Thus, humility is fundamental to the ways of discipleship. The requests in the community prayer evidence this perspective. From daily food to spiritual forgiveness to spiritual protection, the disciple looks to God. There is no moment that can spare dependence on him. 

			God desires prayer that is bold, even shameless, in its approach to him. It is not shameless in the sense of coming to God for all the wants that we have. But it is brash and bold in making use of the access he gives us to seek his face and call on him to develop us spiritually. He desires prayer with nerve. Our requests can be set before his table. The response, of course, is up to God, but the Lord desires we share our heart with him and come with an intense desire to knock at heaven's door. The reason we can be bold emerges in the more subtle argument of the parable: If a human being responds affirmatively to the earthly request of his neighbor, surely a gracious God will respond to our requests about basic spiritual needs. This implicit idea in verses 5--10 emerges explicitly in verses 10--13. We are to come to the Father, knowing that his arms are open for us. I had the gift of a father who loved me. I knew what it meant to be able to walk into his room, his office, the living room, or wherever he was, and know that if I had a crucial need, he had an ear that was leaning in my direction and a heart that was open to my pain. The great lesson of a good father is that it gave me a glimpse of the heavenly Father I now have, even in the absence of my human father. I knew what it was to receive fish or eggs when I asked for them, just as I know now what it is to receive spiritual strength for basic spiritual needs from my Father in heaven. We can be bold because he cares. We can seek his face because he is there, waiting to hear and embrace our spiritual needs.